\subsection*{Проверки воспроизводимости}
Все 71 задание на UMT завершены. Проверены 944 траектории, 998\,400 продуктивных кадров и механический ансамбль каждой серии. Затраты по учёту Slurm составили 920,31 процессорного часа. В основных сериях зонтики гармонически смещают \(Q\), а химическое поле является независимым параметром. Расширение полей при 1243~К позволило получить контакты обеих фаз при двух значениях \(w\). Для более тесной стенки проведены дополнительные серии при 1254, 1262 и 1269~К.

На широкой стенке при 1243~К исходная серия дала \(S=0{,}702\) [0,255; 1,104]. Длинный независимый повтор в трёх полях 3,530, 3,545 и 3,560~эВ дал \(0{,}193\) [\(-0{,}225\); 0,583]. Его диапазон уже интервала между объединёнными полями контактов 3,52767--3,56297~эВ; обе границы лежат снаружи диапазона повтора. Оба корня обнаружены в 35 из 64 повторных блочных выборок этой серии и в 64 из 64 выборок объединённой сетки. После объединения \(S=0{,}494\) [0,092; 0,772], первая и вторая половины дают 0,50 и 0,49.

При 1243~К и 3,545~эВ два старта без зонтика при \(w=0{,}40\) сохранили \(Q\) около 0,79 и 0,11 при практически равных составах 27,21 и 27,19~ат.\% Al. При \(w=0{,}35\) оба старта оказались упорядоченными. Число начальных конфигураций в каждом бассейне задаётся исследователем. Поэтому отношения равновесных весов восстановлены по перекрывающимся зонтикам, а эти старты проверяют поведение фаз независимо.
